\documentclass[10pt,a4paper]{amsart}
\usepackage[margin=19mm]{geometry}
\usepackage{amsmath,amssymb,mathtools}
\usepackage[scr=rsfs,cal=euler]{mathalfa}
\usepackage{xfrac}
\usepackage[unicode,hidelinks,bookmarks=false]{hyperref}
\newcommand{\ie}{i.e.\ }
\newcommand{\cf}{cf.\ }
\newcommand{\resp}{respectively\ }
\newcommand{\Subsection}{\subsection}
\providecommand{\nobreakdash}{\nobreak\hbox{-}}
\setlength{\emergencystretch}{2em}
\multlinegap0pt
\usepackage{amssymb}
\usepackage{mathtools}
\usepackage[shortlabels]{enumitem}
\newtheorem{proposition}{Proposition}
\newtheorem{lemma}{Lemma}
\newtheorem{remark}{Remark}
\newtheorem{theorem}{Theorem}
\title{Schanuel Integration and Euler Characteristic of Semi-algebraic Sets}
\author{Yuxuan Xiao}
\date{Source: arXiv:2607.07180v1 (CC BY 4.0)}
\begin{document}
\maketitle

\noindent\emph{Source and licence.} Schanuel Integration and Euler Characteristic of Semi-algebraic Sets. Version arXiv:2607.07180v1 (CC BY 4.0), distributed by arXiv under the Creative Commons Attribution 4.0 International licence, http://creativecommons.org/licenses/by/4.0/, as recorded in the arXiv licence metadata for this exact version and shown on the abstract page https://arxiv.org/abs/2607.07180v1. Full licence notice: \texttt{licenses/arxiv-2607.07180-CC-BY-4.0.txt}.\par\smallskip

\noindent\textit{Editorial scope.} Counted unit: the article's theorem thm:basis-independence-schanuel-euler (source0.tex lines 574--581). The article's complete original proof of it is delivered with it (source0.tex lines 584--731, one \texttt{proof} environment of 148 lines). No mathematics is written, repaired or re-derived: the statement and the proof are the article's own lines, in the article's own notation, and no retained line is substituted or re-flowed.  Retained uncounted as the source-local context the proof needs: lines 388--393; lines 456--485; lines 529--553; lines 557--572.  Nothing was omitted from any retained span, so the package carries no omission record.  No retained line contains a citation, so the wrapper carries no bibliography and no bibliography entry.  No retained line contains a figure environment, an image-inclusion command or drawing code, so no image is delivered and the package declares no asset dependency.  Editorial formatting only, in addition: an AMS \texttt{amsart} preamble that restates the article's own declarations and macros used by the retained lines, namely the environments \texttt{proposition}, \texttt{lemma}, \texttt{remark}, \texttt{theorem} on separate counters, together with the standard packages those lines need.  The retained environments print in this document as Proposition 1, Lemma 1, Remark 1, Theorem 1 in order of appearance, whereas the article prints the same items with the numbering of its own full document.  The complete native source of the article as distributed in the arXiv e-print for this exact version is bundled unchanged as \texttt{sources/204686/source0.tex}.
\begin{proposition}
\label{prop:semialgebraic-indicator-schanuel-integrable}
Let \(A\subseteq \mathbb{R}^n\) be a semi-algebraic set.  Then, with respect to
any ordered linear basis \((v_1,\ldots,v_n)\) of \(\mathbb{R}^n\), the
indicator function \(\mathbf{1}_A\) is Schanuel-integrable.
\end{proposition}

\begin{lemma}
\label{lem:local-extension-root-function}
Let \(C\subseteq \mathbb{R}^{n-1}\) be connected and locally path-connected,
and let
\[
  \xi:C\longrightarrow \mathbb{R}
\]
be a continuous function such that \(\xi(x)\) is a simple root of a polynomial
\(g\in \mathbb{R}[X_1,\ldots,X_n]\) for every \(x\in C\).  Then for every
\(x_0\in C\), there is an open neighbourhood \(W\subseteq\mathbb{R}^{n-1}\) of
\(x_0\) and a smooth function
\[
  \xi':W\longrightarrow \mathbb{R}
\]
such that
\[
  \xi'|_{W\cap C}=\xi|_{W\cap C}.
\]
Moreover, for \(1\leq k\leq n-1\),
\[
  \frac{\partial \xi'}{\partial X_k}
  =
  -
  \left(
    \frac{\partial g}{\partial X_n}
  \right)^{-1}
  \frac{\partial g}{\partial X_k},
\]
where the right-hand side is evaluated at \((x,\xi'(x))\).
\end{lemma}

\begin{remark}
\label{rem:monotone-root-functions}
Let \(L\subseteq \mathbb{R}[x,y]\) be finite and closed under
\(\partial/\partial y\), and let \(C\subseteq\mathbb{R}\) be a connected cell in
a CAD adapted to \(L\).  Suppose \(\xi:C\to\mathbb{R}\) is a root function
arising from \(L\).  If \(\xi\) has multiplicity \(m\) as a root of
\(f\in L\), then \(\xi\) is a simple root of
\[
  g_\xi
  :=
  \frac{\partial^{m-1}f}{\partial y^{m-1}}.
\]
Thus
\[
  \frac{\partial g_\xi}{\partial y}(x,\xi(x))
  =
  \frac{\partial^m f}{\partial y^m}(x,\xi(x))
  \neq 0
\]
for all \(x\in C\).  If, in addition, the CAD is adapted to all
\(\partial g_\xi/\partial x\), then Lemma~\ref{lem:local-extension-root-function}
implies that the sign of \(\xi'\) is constant on each open interval cell.
Consequently every such root function is either constant or strictly monotone
on each open interval cell.
\end{remark}

\begin{proposition}[Permutation reduction]
\label{prop:permutation-reduction}

Let $f:\mathbb R^n\to\mathbb R$ be a function. Suppose that for every ordered basis
$v=(v_1,\ldots,v_n)$ of $\mathbb R^n$ and every permutation
$\sigma\in S_n$,
\[
\int_{v}f=\int_{v_\sigma}f,
\]
where
\[
v_\sigma=(v_{\sigma(1)},\ldots,v_{\sigma(n)}).
\]
Then $\int_{v}f$ is independent of the choice of the ordered basis $v$.

\end{proposition}

\begin{theorem}
\label{thm:basis-independence-schanuel-euler}
Let \(A\subseteq\mathbb{R}^n\) be semi-algebraic.  Then
\[
  \int_{\varepsilon_{(v_n,\ldots,v_1)}} \mathbf{1}_A
\]
is independent of the ordered basis \((v_n,\ldots,v_1)\).
\end{theorem}

\begin{proof}
By Proposition~\ref{prop:semialgebraic-indicator-schanuel-integrable},
the iterated Schanuel integral is defined for every ordered basis.

Since every ordered basis is obtained from the standard basis by a linear automorphism of $\mathbb R^n$, 
and linear automorphisms preserve semi-algebraic sets, it suffices to consider the standard basis. 
By Proposition~\ref{prop:permutation-reduction}, it is therefore enough to prove that the iterated Schanuel integral is invariant under permutations of the coordinate order.

Every permutation is a product of adjacent transpositions.  Hence it suffices
to prove that
\[
  \int_{\varepsilon} \mathbf{1}_A\,
  dx_n\cdots dx_{i+1}\,dx_i\cdots dx_1
  =
  \int_{\varepsilon} \mathbf{1}_A\,
  dx_n\cdots dx_i\,dx_{i+1}\cdots dx_1
\]
for every adjacent pair \(i,i+1\).  The integrations in the coordinates other
than \(x_i,x_{i+1}\) are the same on both sides and produce a finite
\(\mathbb{Z}\)-linear combination of indicator functions of semi-algebraic
sets.  By finite additivity, the problem reduces to the two-dimensional case:
for every semi-algebraic \(A\subseteq\mathbb{R}^2\), one must show
\[
  \int_{\varepsilon}\mathbf{1}_A\,dx\,dy
  =
  \int_{\varepsilon}\mathbf{1}_A\,dy\,dx .
\]

We prove this two-dimensional statement.  Choose a finite set
\(L\subseteq\mathbb{R}[x,y]\) such that \(A\) is defined by a Boolean
combination of sign conditions on elements of \(L\).  Enlarging \(L\) if
necessary, we may assume that \(L\) is closed under \(\partial/\partial y\).
Now further enlarge \(L\) to a finite set \(L_1\):
$$ L_1= L\cup \left\{\frac{\partial f}{\partial x}: 
\text{if }\frac{\partial f}{\partial x} \text{ is not constant for $f\in L$}\right\} $$
Take a CAD of \(\mathbb{R}^2\)
adapted to \(L_1\).  By monotonicity of the $\text{PROJ}$ construction, this CAD 
refines the CAD adapted to \(L\), and hence \(A\) is a finite disjoint union of
two-dimensional CAD cells for \(L_1\).

It is therefore enough to prove the equality for a single cell
\(K\subseteq\mathbb{R}^2\).  Let \(C=\pi_x(K)\subseteq\mathbb{R}\), where
\(\pi_x(x,y)=x\).  The cell \(K\) is either a graph cell or a band cell over
\(C\).  Let \(d\) be the fiber dimension of \(K\) over \(C\); thus \(d=0\) for a
graph cell and \(d=1\) for a band cell.  By the definition of the Schanuel
integral in the \(y\)-direction,
\[
  \int_{\varepsilon}\mathbf{1}_K\,dy\,dx
  =
  (-1)^d
  \int_{\varepsilon}\mathbf{1}_C\,dx
  =
  (-1)^{d+\dim C}.
\]

We now compute the integral in the opposite order.  There are three cases.

First, suppose \(C\) is a point.  Then \(\dim C=0\).  The projection of \(K\)
to the \(y\)-axis is either a point or an open interval according as \(K\) is a
graph cell or a band cell.  Hence
\[
  \int_{\varepsilon}\mathbf{1}_K\,dx\,dy
  =
  (-1)^d
  =
  \int_{\varepsilon}\mathbf{1}_K\,dy\,dx .
\]

Second, suppose \(C\) is an open interval and \(K\) is a band cell over \(C\).
Write
\[
  K=\{(x,y)\in C\times\mathbb{R}:\xi_1(x)<y<\xi_2(x)\},
\]
allowing \(\xi_1=-\infty\) or \(\xi_2=+\infty\) for unbounded bands.  By
Remark~\ref{rem:monotone-root-functions}, after passing to the refined CAD the
finite root functions are constant or strictly monotone on \(C\).  Therefore
every nonempty horizontal fiber of \(K\) is an open interval.  Since the
projection of \(K\) to the \(y\)-axis is also an interval, we get
\[
  \int_{\varepsilon}\mathbf{1}_K\,dx\,dy
  =
  -\int_{\varepsilon}\mathbf{1}_{\pi_y(K)}\,dy
  =
  (-1)(-1)
  =
  1.
\]
On the other hand,
\[
  \int_{\varepsilon}\mathbf{1}_K\,dy\,dx
  =
  (-1)^{1+1}
  =
  1.
\]

Third, suppose \(C\) is an open interval and \(K\) is the graph of a root
function \(\xi:C\to\mathbb{R}\).  Again, \(\xi\) is either constant or strictly
monotone on \(C\).  If \(\xi\) is strictly monotone, then every nonempty
horizontal fiber is a point and \(\pi_y(K)\) is an open interval.  Hence
\[
  \int_{\varepsilon}\mathbf{1}_K\,dx\,dy
  =
  \int_{\varepsilon}\mathbf{1}_{\pi_y(K)}\,dy
  =
  -1.
\]
If \(\xi\) is constant, then \(\pi_y(K)\) is a point and the horizontal fiber is
the open interval \(C\).  Hence
\[
  \int_{\varepsilon}\mathbf{1}_K\,dx\,dy
  =
  -\int_{\varepsilon}\mathbf{1}_{\pi_y(K)}\,dy
  =
  -1.
\]
In both subcases,
\[
  \int_{\varepsilon}\mathbf{1}_K\,dx\,dy
  =
  -1
  =
  (-1)^{0+1}
  =
  \int_{\varepsilon}\mathbf{1}_K\,dy\,dx .
\]

Thus the two orders of integration agree on every cell \(K\).  Since \(A\) is a
finite disjoint union of such cells and the Schanuel integral is finitely
additive,
\[
  \int_{\varepsilon}\mathbf{1}_A\,dx\,dy
  =
  \sum_K
  \int_{\varepsilon}\mathbf{1}_K\,dx\,dy
  =
  \sum_K
  \int_{\varepsilon}\mathbf{1}_K\,dy\,dx
  =
  \int_{\varepsilon}\mathbf{1}_A\,dy\,dx .
\]
This proves invariance under adjacent transpositions, hence under all
permutations of the coordinate order, and therefore the value of
\[
  \int_{\varepsilon_{(v_n,\ldots,v_1)}} \mathbf{1}_A
\]
is independent of the ordered basis.
\end{proof}

\end{document}
